\section{总结与延伸}

\subsection{一条贯穿三项工作的主线}

Two-phase pretraining 说明数据的\textbf{顺序}有价值；front-loading reasoning 说明能力相关数据的\textbf{时机}有价值；RLP 说明预测前探索的\textbf{行为与奖励}有价值。三者共同把预训练从静态 maximum-likelihood fitting 推向 curriculum-aware、stage-aware、action-aware learning。

\begin{importantbox}{最终知识框架}
\begin{enumerate}
  \item 数据 mixture 必须同时记录 weight 与 temporal schedule。
  \item Quality 与 epoch/repeat 是不同变量，高质量数据也有重复上限。
  \item Reasoning data 的早期分配可以改变后训练可利用的 foundation。
  \item RLP 用 no-think counterfactual 衡量 thought 的 predictive information gain。
  \item Dense reward、EMA baseline 与 group-relative update 共同决定稳定性。
  \item 所有收益都必须在 token、FLOPs、checkpoint、post-training 与 benchmark 口径下解释。
\end{enumerate}
\end{importantbox}

\subsection{从课程结果到研究问题}

第一，能否自动学习连续 curriculum，而不是手工 two-phase？第二，quality classifier 的偏差如何反馈到 model behavior？第三，front-loading 对非 STEM reasoning、multilingual 与 multimodal 是否同样有效？第四，thought reward 能否加入 factuality、causal faithfulness 或 safety 约束而不破坏 scalability？第五，RLP 的 rollout compute 能否通过 latent thought、speculative generation 或更强 batching 降低？

\begin{knowledgebox}{建议的 factorial experiment}
可用 $2\times2\times2$ 设计分别开关 two-phase、front-loading、RLP，在固定 architecture、总 token、近似 FLOPs 与 post-training recipe 下测主效应和交互项。这样才能回答三种策略是互补、冗余还是互相干扰，而不是把三篇论文的结果机械拼接。
\end{knowledgebox}

\subsection{拓展阅读}

建议按问题顺序阅读：先读 Feng 等人的 two-phase pretraining，理解 data blend 与 curriculum；再读 Akter 等人的 front-loading reasoning，检查 reason-base 条件与五条 lesson；最后读 Hatamizadeh 等人的 RLP v2，重点看 expected improvement identity、Algorithm 1、compute matching 与 ablation。随后对照 Quiet-STaR、RPT 与 RLPT，比较各自 reward source 与粒度。

\begin{warningbox}{阅读论文时保留 lecture snapshot}
本讲发生于 2026-04-30：two-phase 使用 arXiv v1，front-loading 使用 v1，RLP 使用 2026-03-01 的 v2。后续 revision、复现或产品结果可以作为新证据，但不能倒推成课堂当日已经成立的事实。
\end{warningbox}

最终，课程提出的“next-generation intelligence”不是一个单独的新模型，而是一种训练观：在有限数据与算力下，模型应更聪明地选择学习顺序、更早建立 reasoning foundation，并让探索行为得到可扩展、可比较的反馈。预训练的未来，可能不只在更大的 $N$，也在更好的 learning process。

\end{document}
